% F6 -- Four runs, four datasets, the same shape: a plateau under the baseline.
\begin{figure}[t]
\centering
\plotlink{\nbllm}{\fitwidth{%
\begin{tikzpicture}
\begin{groupplot}[
  group style={group size=2 by 2, horizontal sep=1.6cm, vertical sep=2.0cm,
               xlabels at=edge bottom, ylabels at=edge left},
  rep, width=0.46\linewidth, height=4.3cm,
  xlabel={epoch}, ylabel={validation top-1 (\%)},
  ymin=0, ymax=100, ytick={0,20,...,100},
  no markers, title style={yshift=2pt},
]
% ---------- (a) D1, unfiltered day, 500 users, 49 epochs ----------
\nextgroupplot[title={\footnotesize (a) \sraw, 49 epochs},
               xmin=0, xmax=50, xtick={1,10,20,30,40,49}]
\addplot[clbase,dashed,thick,domain=0:50,samples=2]{69.2};
\addplot[clmodel,thick,mark=*,mark size=0.8pt] coordinates {
(1,18.9)(2,53.8)(3,55.8)(4,61.5)(5,56.4)(6,63.9)(7,62.7)(8,60.3)(9,65.0)(10,61.8)
(11,60.8)(12,59.6)(13,59.8)(14,61.9)(15,58.7)(16,59.8)(17,60.6)(18,61.1)(19,60.2)(20,55.8)
(21,53.3)(22,53.5)(23,53.0)(24,50.7)(25,50.2)(26,50.5)(27,50.7)(28,47.9)(29,49.0)(30,47.8)
(31,46.6)(32,47.0)(33,48.0)(34,48.1)(35,47.2)(36,46.7)(37,48.3)(38,49.0)(39,47.6)(40,46.9)
(41,47.4)(42,46.9)(43,46.9)(44,47.9)(45,47.4)(46,47.4)(47,47.9)(48,48.1)(49,48.1)};
\node[star,star points=5,fill=clgain,draw=none,inner sep=1.6pt] at (axis cs:9,65.0) {};
\node[font=\scriptsize,anchor=south west,text=clgain] at (axis cs:9,66) {peak, epoch 9};
\node[font=\scriptsize,anchor=south east,text=clbase] at (axis cs:49,70) {baseline 69.2\,\%};
%
% ---------- (b) D2, filtered, 14 epochs ----------
\nextgroupplot[title={\footnotesize (b) \sfilt, 14 epochs},
               xmin=0, xmax=15, xtick={1,4,...,14}]
\addplot[clbase,dashed,thick,domain=0:15,samples=2]{79.0};
\addplot[clmodel,thick,mark=*,mark size=1pt] coordinates {
(1,56.1)(2,72.2)(3,73.2)(4,73.5)(5,71.8)(6,73.8)(7,71.8)(8,71.8)(9,69.8)(10,70.8)
(11,69.1)(12,67.6)(13,68.2)(14,67.2)};
\node[star,star points=5,fill=clgain,draw=none,inner sep=1.6pt] at (axis cs:6,73.8) {};
\node[font=\scriptsize,anchor=south,text=clgain] at (axis cs:6,75) {epoch 6};
\node[font=\scriptsize,anchor=south east,text=clbase] at (axis cs:14.5,80) {79.0\,\%};
%
% ---------- (c) D3, mobile users, 9 epochs ----------
\nextgroupplot[title={\footnotesize (c) \smob, 9 epochs},
               xmin=0, xmax=10, xtick={1,3,...,9}]
\addplot[clbase,dashed,thick,domain=0:10,samples=2]{43.5};
\addplot[clmodel,thick,mark=*,mark size=1pt] coordinates {
(1,19.8)(2,40.2)(3,42.3)(4,42.5)(5,43.0)(6,42.3)(7,39.9)(8,38.7)(9,38.4)};
\node[star,star points=5,fill=clgain,draw=none,inner sep=1.6pt] at (axis cs:5,43.0) {};
\node[font=\scriptsize,anchor=south,text=clgain] at (axis cs:5,44.5) {epoch 5};
\node[font=\scriptsize,anchor=north east,text=clbase] at (axis cs:9.7,41.5) {43.5\,\%};
%
% ---------- (d) D3 + hour tokens, 11 epochs ----------
\nextgroupplot[title={\footnotesize (d) \smob + hour tokens, 11 epochs},
               xmin=0, xmax=12, xtick={1,3,...,11}]
\addplot[clbase,dashed,thick,domain=0:12,samples=2]{45.5};
\addplot[clmodel,thick,mark=*,mark size=1pt] coordinates {
(1,21.1)(2,39.0)(3,41.9)(4,41.1)(5,42.2)(6,41.3)(7,40.0)(8,39.4)(9,38.1)(10,38.8)(11,37.1)};
\node[star,star points=5,fill=clgain,draw=none,inner sep=1.6pt] at (axis cs:5,42.2) {};
\node[font=\scriptsize,anchor=south,text=clgain] at (axis cs:5,43.6) {epoch 5};
\node[font=\scriptsize,anchor=north east,text=clbase] at (axis cs:11.7,43.5) {45.5\,\%};
\end{groupplot}
\end{tikzpicture}}}
\caption{Validation accuracy per epoch on four datasets}
\label{fig:validation-curves}
\figtext{%
The same shape, four times, on four different datasets. In each panel the
horizontal axis is the training epoch and the vertical axis is top-1 accuracy
on validation users the run never trains on; the solid line is the model, the
star its best epoch, the dashed line that dataset's persistence baseline.
Validation climbs for four to nine epochs, peaks, then decays for tens of
epochs while the training loss keeps falling. In panel~(a) the peak is
65.0\,\% at epoch~9 against 48.1\,\% at the last epoch, so reading the end of
a run instead of its best point costs sixteen points and forty epochs of GPU
time. Note that the dashed line moves between panels, 69.2\,\% on the
unfiltered day and 43.5\,\% once the corpus is restricted to mobile users,
because the users themselves are different; the panels are comparable only on
the distance between the two lines, and in all four that distance is negative.
Per-epoch values are recovered by digitising each run's own validation curve,
with the vertical axis calibrated on the two values printed on that curve, so
an individual point carries about $\pm1\pt$ of reading error; the figures
quoted in the text are the reported ones. Qwen2.5-0.5B + QLoRA throughout.}
\figsource{\nbllm}{train\_cellid\_llm.ipynb}{the training loop whose per-epoch validation scores these curves are}
\end{figure}
